\documentclass[10pt,a4paper]{amsart}
\usepackage[margin=19mm]{geometry}
\usepackage{amsmath,amssymb,mathtools}
\usepackage[scr=rsfs,cal=euler]{mathalfa}
\usepackage{xfrac}
\usepackage[unicode,hidelinks,bookmarks=false]{hyperref}
\newcommand{\ie}{i.e.\ }
\newcommand{\cf}{cf.\ }
\newcommand{\resp}{respectively\ }
\newcommand{\Subsection}{\subsection}
\providecommand{\nobreakdash}{\nobreak\hbox{-}}
\setlength{\emergencystretch}{2em}
\multlinegap0pt
\usepackage{bm}
\usepackage{bbm}
\usepackage{dsfont}
\usepackage{xspace}
\usepackage{cancel}
\usepackage{mathtools}
\usepackage{amsthm}
\makeatletter
\@ifundefined{theorem}{\newtheorem{theorem}{Theorem}}{}
\@ifundefined{claim}{\newtheorem{claim}[theorem]{Claim}}{}
\@ifundefined{lemma}{\newtheorem{lemma}[theorem]{Lemma}}{}
\makeatother
\title[A universal threshold for geometric embeddings of trees]{A universal threshold for geometric embeddings of trees}
\author{Dylan J. Altschuler, Pandelis Dodos, Konstantin Tikhomirov, Konstantinos Tyros}
\date{Source: arXiv:2504.15212v2 (CC BY 4.0)}
\begin{document}
\maketitle

\noindent\emph{Source and licence.} A universal threshold for geometric embeddings of trees. Version arXiv:2504.15212v2 (CC BY 4.0), distributed under the Creative Commons Attribution 4.0 International licence, as recorded in the arXiv licence metadata for this exact version and shown on the abstract page https://arxiv.org/abs/2504.15212v2. Full rights notice: licenses/arxiv-2504.15212-CC-BY-4.0.txt.\par\smallskip

\noindent\textit{Editorial scope.} Counted unit: the Claim whose source label is \texttt{claim\_Lovatz\_condition} in the article (source0.tex lines 544--546, source label \texttt{claim\_Lovatz\_condition}), delivered together with the article's own complete original proof of it (source0.tex lines 547--613, one \texttt{proof} environment of 67 lines). No mathematics is written, repaired or re-derived: statement and proof are the article's own text line for line. Required context retained uncounted: thm:thinshell (lines 248--253); thm:slicing (lines 257--262); lem:ball\_volume (lines 316--318); lem:prob\_bound\_1 (lines 330--340); lem:prob\_bound\_2 (lines 356--363); eq:001 (lines 411--413); eq:009 (lines 450--453). Every definition, display and label of those lines is retained unchanged. Omitted from this package and not redistributed: 1--247, 254--256, 263--315, 319--329, 341--355, 364--410, 414--449, 454--543, 614--912. Nothing inside a retained excerpt is omitted or altered: every retained body is delivered exactly as the article's lines are written, so no omission receipt of any kind accompanies this package. Citations: the retained excerpts carry no citation command at all, so no bibliography entry of the article is transcribed and no reference is redistributed. Reference adaptation: none was needed and none was made; every reference inside the delivered unit resolves inside it, so no unresolved reference or citation is produced. Printed slips kept literal: no printed word, symbol, hypothesis or number of the counted statement or of its proof was altered. Layout and class: the article's own class is \texttt{amsart} with packages including graphicx, xcolor, fullpage, amsmath, amsthm, amssymb, latexsym, color, hyperref, lipsum, enumitem, bbm, tikz, subfigure; the wrapper is the lane's pinned \texttt{amsart} editorial wrapper at 10pt on A4, which restates the theorem-like environments the retained text needs and adds the article's own macro definitions that the retained text needs (none), with extra wrapper packages (bm, bbm, dsfont, xspace, cancel, mathtools). The delivered numbering is the unit's own. Assets: no retained span contains a figure environment, a graphics-inclusion command, a PDF-page inclusion, a table or a TikZ picture; no image file is copied into this package, the package declares no asset dependency and no image file is redistributed with it. Licence: version arXiv:2504.15212v2 of this article is distributed under the Creative Commons Attribution 4.0 International licence, as recorded in the arXiv licence metadata for this exact version and shown on the abstract page https://arxiv.org/abs/2504.15212v2. A full rights notice is delivered at licenses/arxiv-2504.15212-CC-BY-4.0.txt. Characters: every retained excerpt is pure ASCII, so the delivered engine, XeLaTeX with its default Latin Modern text font, reports no missing character.
\begin{theorem}[Thin-shell concentration for log-concave measures] \label{thm:thinshell}
There exist absolute constants $0<\kappa <1$ and $C>0$ with the following property. If\, $\mathbf{U}$ is an isotropic random vector in $\mathbb{R}^m$ with a log-concave\footnote{A function $f\colon \mathbb{R}^m\to [0,\infty)$ is \emph{log-concave} if, setting $K:=\{x\in\mathbb{R}^m\colon f(x) > 0\}$, we have that the function $\log f\colon K\to\mathbb{R}$ is concave. (Note that the set $K$ is convex.)} distribution, then
\begin{equation} \label{eq2.2}
\mathbb{P}\Big(\big|\|\mathbf{U}\|_2 - \sqrt{m}\big|  \geqslant m^{\frac12 - \kappa} \Big) \leqslant C \exp\big(-m^{\kappa}\big).
\end{equation}
\end{theorem}

\begin{theorem}[Anticoncentration for log-concave measures via slicing] \label{thm:slicing}
There exists an absolute constant $L>0$ with the following property. If\, $\mathbf{U}$ is an isotropic random vector in $\mathbb{R}^m$ with a log-concave density $f$, then for any $x \in \mathbb{R}^m$,
\begin{equation} \label{eq2.3}
f(x) \leqslant L^{m}.
\end{equation}
\end{theorem}

\begin{lemma} \label{lem:ball_volume}
There exists an absolute constant $C_B>0$ such that for any $m$-dimensional normed space $X=(\mathbb{R}^m,\|\cdot\|_X)$ in isotropic position, we have $\lambda(B_X)\leqslant C_B^m$, where $\lambda$ denotes the Lebesgue measure in $\mathbb{R}^m$.
\end{lemma}

\begin{lemma} \label{lem:prob_bound_1}
Let $X=(\mathbb{R}^m,\|\cdot\|_X)$ be a normed space in isotropic position, let $k$ be a positive integer, and let $0<c_1<\frac{1}{2}$ such that
\begin{equation} \label{eq:004}
k\geqslant \big(C_B\,L\big)^{\frac{4}{1-2c_1}},
\end{equation}
where $C_B$ is as in Lemma \ref{lem:ball_volume}, and $L$ is as in Theorem \ref{thm:slicing}. Finally, let $(\mathbf{Y}_1,\dots,\mathbf{Y}_k)$ be i.i.d. random vectors uniformly distributed in $B_X$. Then,
\begin{equation} \label{eq3.1}
\mathbb{P} \Bigg( \Big\|  \sum_{i=1}^{k} \mathbf{Y}_i \Big\|_X \leqslant k^{c_1} \Bigg)
\leqslant \exp\Big(-\frac{1}{2} \Big(\frac{1}{2}-c_1\Big) m \log k \Bigg).
\end{equation}
\end{lemma}

\begin{lemma} \label{lem:prob_bound_2}
There exists a positive integer $m_1$ with the following property. Let $m\geqslant m_1$ be an integer, let $X=(\mathbb{R}^m,\|\cdot\|_X)$ be a normed space in isotropic position, let $k\geqslant 2$ be an integer, and let $(\mathbf{Y}_1,\dots,\mathbf{Y}_k)$ be i.i.d. random vectors uniformly distributed in $B_X$. Then,
\begin{equation} \label{eq3.4}
\mathbb{P} \Bigg( \Big\|  \sum_{i=1}^{k} \mathbf{Y}_i \Big\|_X \leqslant 1+\frac{1}{2\,m^{1-\kappa}} \Bigg)
\leqslant \exp\big(-m^{\frac{\kappa}{2}}\big),
\end{equation}
where $\kappa$ is as in Theorem \ref{thm:thinshell}.
\end{lemma}

\begin{equation} \label{eq:001}
\alpha_1:=64, \ \ \ \ c_1:=\frac{1}{4}, \ \ \ \ c_2:=\min\Big\{\frac{1}{4},\frac{\kappa}{4}\Big\},
\end{equation}

\begin{equation} \label{eq:009}
\ell_0 := \exp\big( \log^{c_2}N\big) \ \ \ \text{ and } \ \ \
k_0:= \max\Big\{ 3,\big\lceil (C_BL)^{4/(1-2c_1)}\big\rceil, \big\lceil 4^{1/c_1}\big\rceil \Big\},
\end{equation}

\begin{claim}\label{claim_Lovatz_condition}
For every $\{u,v\}\in\mathcal{V}$, we have $\mathbb{P}(\mathcal{L}_{u,v})\leqslant\frac{1}{2}h(\{u,v\})$.
\end{claim}

\begin{proof}
Fix $\{u,v\}\in\mathcal{V}$, and set $k:=\mathrm{dist}_T(u,v)$. We distinguish three cases.

First, assume that $k\in \{2,\dots,k_0\}$. Since the random vectors
$\sum_{e\in P(r_T,u)}\mathbf{Y}_e - \sum_{e\in P(r_T,v)}\mathbf{Y}_e$ and
$\sum_{e\in P(u,v)}\mathbf{Y}_e $ have the same distribution, then provided that $N_0$ is sufficiently large,
by Lemma \ref{lem:prob_bound_2},
\begin{equation}\label{eq:019}
\mathbb{P}(\mathcal{L}_{u,v})\leqslant \exp\big(- m ^{\frac{\kappa}{2}}\big) \leqslant
\exp\Big(- \alpha_1^{\frac{\kappa}{2}}\frac{\log^{\frac{\kappa}{2}}N}{(\log\log N)^{\frac{\kappa}{2}}}\Big).
\end{equation}
On the other hand, by the choice of $\ell_0$ in \eqref{eq:009} and the fact that $b_k\geqslant b_{k_0}$ for $k\in \{2,\dots, k_0\}$,
\begin{equation}\label{eq:020}
\frac{h(\{u,v\})}{2} =\frac{b_k}{2} \geqslant \frac{b_{k_0}}{2} \geqslant
\exp\big(-\log 24-k_0\log(2\Delta-2)-\log k_0 - \log^{c_2}N\big).
\end{equation}
Since $c_2<\frac{\kappa}{2}$ and using the fact that $N_0$ can be taken sufficiently large, the claim
follows in this case by~\eqref{eq:019} and \eqref{eq:020}.

Next, assume that $k\in \{k_0+1,\dots,k_1\}$. Using again the fact that
$\sum_{e\in P(r_T,u)}\mathbf{Y}_e - \sum_{e\in P(r_T,v)}\mathbf{Y}_e$ and
$\sum_{e\in P(u,v)}\mathbf{Y}_e $ have the same distribution, by Lemma \ref{lem:prob_bound_1},
the choice of $k_0$ in \eqref{eq:009} and the fact that $c_1\leqslant\frac{1}{4}$, we obtain that
\begin{equation}\label{eq:021}
\mathbb{P}(\mathcal{L}_{u,v}) \leqslant \exp\Big(-\frac{1}{2} \Big(\frac{1}{2}-c_1\Big)m\log k \Big)
\leqslant \exp\Big( -\frac{\alpha_1}{8} \cdot \frac{\log N}{\log\log N} \log k \Big).
\end{equation}
Since the function $f(x) = \frac{x}{\log x}$ is increasing for $x\geqslant3$ and $k>k_0\geqslant3$, we see that
\begin{align*}
\frac{h(\{u,v\})}{2} & =
\exp\big(-\log 24-k\log(2\Delta-2)-\log k - \log^{c_2}N\big) \\
& =\exp\Big(-\Big(\frac{\log 24}{\log k}+\frac{k}{\log k }\log(2\Delta-2)+1+
\frac{\log^{c_2}N}{\log k}\Big)  \log k \Big)\\
& \geqslant \exp\Big(-\Big(\frac{\log 24}{\log 3 }+\frac{k_1}{\log k_1 }\log(2\Delta-2)+1+
\frac{\log^{c_2}N}{\log k }\Big)  \log k \Big)\\
& \geqslant \exp\Big(-\Big(4+\frac{\log N }{\log\log N-\log\log(\Delta-1)}\cdot\frac{\log (2(\Delta-1))}{\log(\Delta-1)}+
\frac{\log^{c_2}N}{\log k}\Big)  \log k \Big).
\end{align*}
Taking $N_0$ sufficiently large and observing that $\frac{\log(2(\Delta-1))}{\log(\Delta-1)}\leqslant2$
and $\frac{\alpha_1}{8}>4$, we conclude that
\begin{equation}\label{eq:022}
\frac{h(\{u,v\})}{2}
\geqslant\exp\Big(-\Big(4+4\frac{\log N}{\log\log N}+\log^{c_2}N\Big)\log k\Big)
\stackrel{\eqref{eq:021}}{\geqslant}\mathbb{P}(\mathcal{L}_{u,v}).
\end{equation}

Finally, assume that $k\in \{k_1+1,\dots,\ell_0\}$. As before, since
$\sum_{e\in P(r_T,u)}\mathbf{Y}_e - \sum_{e\in P(r_T,v)}\mathbf{Y}_e$ and
$\sum_{e\in P(u,v)}\mathbf{Y}_e $ have the same distribution, by Lemma \ref{lem:prob_bound_1},
the choice of $k_0$ in \eqref{eq:009} and the fact that $k>k_0$ and $c_1\leqslant\frac{1}{4}$, we have
\begin{equation}\label{eq:023}
\mathbb{P}(\mathcal{L}_{u,v}) \leqslant \exp\Big(-\frac{1}{2} \Big(\frac{1}{2}-c_1\Big)m\log k \Big)
\leqslant \exp\Big( -\frac{\alpha_1}{8} \cdot \frac{\log N}{\log\log N} \log k \Big).
\end{equation}
Taking $N_0$ sufficiently large, we may assume that
$\log k \geqslant \log \log N - \log\log(\Delta-1) \geqslant \frac{1}{2}\log\log N$.
Thus, by \eqref{eq:023} and the choice of $\alpha_1$ in \eqref{eq:001}, we have that
\begin{equation}\label{eq:024}
\mathbb{P}(\mathcal{L}_{u,v})\leqslant \exp( -4\log N).
\end{equation}
Since $k_1 < k \leqslant \ell_0$, taking $N_0$ sufficiently large, we see that
\begin{equation} \label{eq:025}
\frac{h(\{u,v\})}{2} = \exp\Big(-\log\Big(\frac{4\pi^2}{3}\Big) -2\log k -2\log N\Big)
\geqslant \exp\Big(- 3\log N \Big).
\end{equation}
Thus, this final case the claim follows from  \eqref{eq:024} and \eqref{eq:025}.
\end{proof}

\end{document}
